\bigskip
\begin{center}\textls[120]{e) Die Inauguraldissertation.}\end{center}
\HistoricalFoot{X\,{\scriptsize 2} Abh. 2.}{8}
\indent In der Inauguraldissertation\footnote{\textit{Demonstratio nova theorematis omnem functionem algebraicam rationalem integram unius variabilis in factores reales primi vel secundi gradus resolvi posse}, Helmstedt, Fleckeisen 1799, Werke III, S. 1 ff.,
deutsch in \textsc{Ostwalds} Klassikern Nr. 14.} nimmt \textsc{Gauss} in bezug auf die komplexen
oder, wie er sie daselbst noch bezeichnet, imaginären Grössen den Standpunkt
ein, dass er erklärt, ihre Zulässigkeit sei bis jetzt noch nicht in ausreichender
Weise nachgewiesen worden, und er wolle deshalb von ihrem Gebrauche absehen,
indem er sich jedoch »die Rechtfertigung ihrer Einführung sowie eine
eingehende Erörterung dieser ganzen Sache für eine andere Gelegenheit vorbehält«
(Werke III, S. 6, Fussnote). Die Dissertation selbst zerfällt in zwei\par
\clearpage\SourceLeaf{143}\ArithmeticHead{58}{L}{L. SCHLESINGER, ÜBER GAUSS’ FUNKTIONENTHEORETISCHE ARBEITEN.}
\noindent Teile, einen kritischen, in dem \textsc{Gauss} die Beweise, die \textsc{d’Alembert}, \textsc{Euler},
\textsc{de Foncenex} und \textsc{Lagrange} für die Wurzelexistenz gegeben haben, bespricht,
und einen zweiten, in dem er seinen eigenen Beweis entwickelt. — Aus dem
ersten, »einer unerschöpflichen Fundgrube mathematischer Ideen«\footnote{\textsc{Frobenius}, \textit{Gedächtnisrede auf L. Kronecker}, Abhandlungen der Berliner Akademie 1893, S. 10.}, wollen
wir nur einige für unseren Gegenstand bedeutungsvolle Momente hervorheben.
— In der Kritik des \textsc{d’Alembertschen} Beweises (Art. 6) enthält die Fussnote
zum dritten Einwand (Werke III, S. 10) die schon erwähnten Bemerkungen
über Konvergenz und Divergenz von Potenzreihen. Es gebe Potenzreihen, die
stets divergieren, wie klein auch der Wert der Variabeln genommen werde,
z. B. die hypergeometrischen, das sind solche, in denen der Koeffizient des
\(n\)-ten Gliedes das Produkt der \(n\) ersten Glieder einer arithmetischen Reihe
ist, aber unter den stets divergenten Potenzreihen gebe es solche, die zuerst
sehr stark, dann immer schwächer und schwächer konvergieren (d. h. dass
die Glieder immer kleiner werden) und schliesslich immer mehr divergieren (d. h.
dass die Glieder wachsen), die aber trotzdem die Summe (was darunter zu
verstehen ist, wird nicht weiter erklärt) mit grosser Genauigkeit liefern, wenn
nicht zu viele Glieder genommen werden. Es sei erwünscht, klar und präzis
zu zeigen, warum dies stattfindet und wie man den Fehler abschätzen könne.
— \textsc{Gauss} hat sich also damals für divergente Reihen von asymptotischem
Charakter interessiert, wie wir bereits oben bemerkt haben. — Besonders
wichtig ist aber der vierte Einwand, der besagt, dass eine veränderliche Grösse
\(\mathrm{X}\), die eine (kleinste) obere Schranke besitzt, diese nicht zu erreichen braucht,
obgleich sie ihr beliebig nahe kommen kann. \textsc{Gauss} bemerkt zwar, dass in
dem Falle, wo \(\mathrm{X}\) eine ganze rationale Funktion bedeutet, dies nicht eintritt,
und fügt hinzu, dass der wahre Nerv des \textsc{d’Alembertschen} Beweises von allen
Einwürfen unberührt bleibe, so dass man auf ihn, mit grösserer Sorgfalt als
bei \textsc{d’Alembert}, nicht nur einen strengen Beweis des algebraischen Satzes
aufbauen, sondern aus ihm auch alles ablesen könnte, was sich in bezug auf
transzendente Gleichungen wünschen lässt. In der Tat deutet \textsc{Gauss} im Art. 24
eine solche Ausführung des \textsc{d’Alembertschen} Beweises in geometrischer Form
an, und \textsc{Clausen} hat später (Astronomische Nachrichten 17, 1840, Spalte 325 ff.)
eine analytische Einkleidung desselben Prinzips gegeben. Freilich unterliegt
\textsc{Clausens} Beweis, wie \textsc{Gauss} in einem Briefe an \textsc{Schumacher} (Werke X, 1,\par
\clearpage\SourceLeaf{144}\ArithmeticHead{59}{R}{III e. DIE INAUGURALDISSERTATION.}
\HistoricalFoot{}{\(8^*\)}
\noindent S. 108) betont, dem vierten Einwand, und der Beweis, dass dieser sich für
eine ganze rationale Funktion beheben lässt, kann, wie \textsc{Gauss} daselbst sagt,
»im Geiste des Rigor antiquus nicht ohne einige Umständlichkeit geschehen«.

Ebenso wie beim Beweise von \textsc{d’Alembert} ist \textsc{Gauss’} Kritik auch bei
den Beweisen von \textsc{Euler} und \textsc{de Foncenex} nicht nur negativ, sondern auch
aufbauend. Den Grundgedanken des ersten Beweises von \textsc{Euler} (Art. 7, 8
der Dissertation) und des von \textsc{de Foncenex} (Art. 10, 11) hat \textsc{Gauss} in seinem
zweiten Beweise (1815, Werke III, S. 31) in exakter Weise durchgeführt, und
in bezug auf den zweiten \textsc{Eulerschen} Beweis (Art. 9), der von der Annahme
ausgeht, dass jede algebraische Gleichung durch »arithmetische Operationen
und Wurzelausziehungen« gelöst werden kann, sagt er, es möchte vielleicht
nicht allzuschwierig sein, die Unmöglichkeit einer solchen Lösung schon für
den fünften Grad in aller Strenge nachzuweisen. Ähnlich spricht er sich,
unter Bezugnahme auf diese Stelle der Dissertation, im Art. 359 der \textit{Disqu.
arithm.} (Werke I, S. 449) aus, während er noch in der Tagebuchnotiz 37. vom
September 1796 über einen Versuch berichtet, der zur Lösung einer allgemeinen
Gleichung beliebigen Grades durch Wurzelzeichen mit Hilfe der \textsc{Lagrangeschen}
Resolvente führen sollte (siehe Werke X, 1, S. 503, 504 und die
darauf bezügliche Anmerkung von \textsc{Loewy}).

In bezug auf den vierten Einwand gegen \textsc{d’Alembert} möge noch hervorgehoben
werden, dass die darin betonte Notwendigkeit, zu unterscheiden, ob
eine veränderliche Grösse ihrer oberen Schranke nur beliebig nahe kommt
oder ob sie diese wirklich erreicht, in \textsc{Gauss’} eigener wissenschaftlicher Produktion
eine merkwürdige Rolle spielt. In seinem dritten Beweise für die
Wurzelexistenz (1816, Werke III, S. 57), auf den wir noch ausführlich zurückzukommen
haben, hat er zu Beginn des Art. 2 gegen diesen Einwand verstossen,
und ebenso hat er in der Abhandlung über die im verkehrten Verhältnis
des Quadrats der Entfernung wirkenden Kräfte (1840, Werke V, S. 197)
in den Artikeln 30, 31 aus dem Vorhandensein einer unteren Schranke für
ein gewisses Integral ohne weiteres gefolgert, dass dieses Integral einen kleinsten
Wert wirklich annimmt (vergl. \textsc{Bolza}, Abhandl. 5 dieses Bandes, S. 77, 78).
Auch lobt er in dem Briefe an \textsc{Drobisch} (1834, Werke X, 1, S. 106) neben
seinem eigenen dritten Wurzelexistenzbeweise den von \textsc{Cauchy} in der \textit{Analyse
algébrique} (1821, S. 329) gegebenen, obwohl dieser genau wie die Beweise von\par
\clearpage\SourceLeaf{145}\ArithmeticHead{60}{L}{L. SCHLESINGER, ÜBER GAUSS’ FUNKTIONENTHEORETISCHE ARBEITEN.}
\noindent \textsc{d’Alembert} und \textsc{Clausen} dem vierten Einwand unterliegt (vergl. die Bemerkung
zu dem Briefe an \textsc{Drobisch}, Werke X, 1, S. 110), und auch in der Einleitung
zur Jubiläumsschrift\footnote{\textit{Beiträge zur Theorie der algebraischen Gleichungen} 1849, Werke III, S. 71 ff.} erwähnt er diesen \textsc{Cauchyschen} Beweis neben seinen
drei eigenen.

Der zweite Teil der Dissertation (Artikel 13—24) wird durch die geflissentliche
Umgehung des Gebrauchs komplexer Grössen etwas weniger durchsichtig
als die Jubiläumsschrift, in der \textsc{Gauss} bei seinem fünfzigjährigen Doktorjubiläum
(1849) eine neue Darstellung seines alten Beweises von 1799 mit
Benutzung komplexer Grössen gab, da, wie er sagt, »gegenwärtig .\ . der Begriff
der komplexen Grössen jedermann geläufig« sei. — Setzt man in
\[
\mathrm{X}=x^n+Ax^{n-1}+\ldots+M
\]
\(x=t+iu,\ t=r\cos\varphi,\ u=r\sin\varphi\), so wird \(\mathrm{X}=T+iU\), und \textsc{Gauss} deutet
nun (Art. 16) die \(r,\varphi\) als Polarkoordinaten in einer Ebene. — Dies sind
offenbar die »Spuren« der geometrischen Deutung komplexer Grössen, die »der
aufmerksame Leser« in der Dissertation »wiederfindet«\footnote{Vergl. oben S. 53. Wie \textsc{P. Stäckel}, a. a. O. \textsuperscript{67)}, S. 124 hervorhebt, hat \textsc{d’Alembert} in seiner
Arbeit \textit{Recherches sur le calcul intégral}, Histoire de l’Académie 1746 (Berlin, 1748), S. 182—200, derselben,
die \textsc{Gauss} in den Artikeln 5 und 6 seiner Dissertation (Werke III, S. 7) kritisch bespricht, die Gleichungen
\(T=0,\ U=0\) als Kurven in der \((t,u)\)-Ebene aufgefasst, die durch die rechtwinkligen Koordinaten ihrer
Schnittpunkte diejenigen Werte \(t,u\) bestimmen, für die \(t+u\sqrt{-1}\) eine Wurzel der Gleichung \(\mathrm{X}=0\)
darstellt. Vergl. ausser der von \textsc{Stäckel} zitierten Stelle (Histoire, S. 191, Remarque I-ère) besonders auch
a. a. O., S. 188. Immerhin ist es noch ein bedeutsamer Schritt, von dieser Deutung der \textls[120]{Wurzeln} \(t+ui\)
bis zu der Darstellung einer \textls[120]{beliebigen} komplexen Grösse \(t+ui\) durch die Punkte einer Ebene.}. \textsc{Gauss} sagt nun
(Art. 15), dass die Existenz eines reellen Faktors ersten oder zweiten Grades
von \(\mathrm{X}\) bewiesen sei, wenn man zeigt, dass sich \(r,\varphi\) so bestimmen lassen,
dass die Gleichungen \(T=0,\ U=0\) gleichzeitig erfüllt sind, und er erschliesst
aus dem Verlaufe der durch diese Gleichungen gegebenen algebraischen Kurven
für grosse Werte von \(r\), dass Schnittpunkte dieser Kurven vorhanden sein
müssen. Der Art. 18 zeigt auch deutlich, dass dieser in seiner Einfachheit
so bewundernswerte Gedanke, \textsc{Gauss’} Interesse an dem asymptotischen Verhalten
von Funktionen seine Entstehung verdankt. In der Tat ist \(\lim\limits_{r\to\infty}\mathrm{X}\,.\,x^{-n}=1\), d. h.
es wird \(\mathrm{X}\) für grosse Werte von \(r\) durch \(x^n\) asymptotisch dargestellt, so dass,
wie es im Art. 18 heisst, die Kurven \(T=0,\ U=0\) die in je \(2n\) Strahlen zerfallenden
Kurven \(\cos n\varphi=0\) bezw. \(\sin n\varphi=0\) zu Asymptoten haben. Dem-\par
\clearpage\SourceLeaf{146}\ArithmeticHead{61}{R}{III e. DIE INAUGURALDISSERTATION.}
\noindent gemäss besitzen die Kurven \(T=0,\ U=0\) je \(2n\) sich ins Unendliche erstreckende
Zweige, die auf einem Kreise mit sehr grossem Radius je \(2n\)
asymptotisch äquidistante Punkte ausschneiden, so zwar, dass beim Durchlaufen
der Kreisperipherie in einem bestimmten Sinne immer abwechselnd ein Schnittpunkt
von \(T=0\) und einer von \(U=0\) getroffen wird. In den Artikeln 19 und
20 wird dies für einen Kreis, ausserhalb dessen sicher kein Schnittpunkt der
Kurven \(T=0,\ U=0\) liegt, noch direkt gezeigt, und im einzelnen verfolgt,
wie beim Durchlaufen der Peripherie die Vorzeichen der \(T,U\) aufeinander
folgen. Daraus ergibt sich anschaulich unmittelbar, dass die beiden Kurven
\(T=0,\ U=0\) sich innerhalb dieses Kreises notwendig schneiden müssen, ja
man kann, wie es im Art. 23 heisst, ebenso zeigen, dass es im allgemeinen
\(n\) Schnittpunkte geben wird. — In eben diesem Artikel sagt \textsc{Gauss} auch,
dass es durchaus nicht unmöglich sei, den von ihm auf geometrische Prinzipien
begründeten Beweis in rein analytischer Form zu geben, er glaube
aber, dass der wahre Nerv des Beweises bei der geometrischen Einkleidung
viel klarer vor Augen treten würde, als es sich bei einem analytischen Beweise
erwarten liesse. Eine solche analytische Fassung (Arithmetisierung) des
\textsc{Gaussschen} Beweises hat \textsc{Ostrowski} gegeben\begingroup\renewcommand{\thefootnote}{79a)}\footnote{\textsc{A. Ostrowski}, Anhang zum Heft VIII der \textit{Materialien} 1920, S. 50 ff. (Abh. 3 dieses Bandes).}\endgroup\addtocounter{footnote}{-1}.

Man kann dieser Schlussweise auch die folgende Wendung geben: Es mögen
die Koeffizienten von \(\mathrm{X}\) reelle Grössen sein, wie \textsc{Gauss} in der Dissertation
stets voraussetzt, was ja übrigens keine erhebliche Beschränkung darstellt;
nehmen wir dann auf der Kreisperipherie mit grossem Radius etwa den Punkt
zum Ausgangspunkt, in dem der Kreis von der positiven reellen Achse geschnitten
wird, so ist dort \(U=0\), also \(\mathrm{X}\) reell positiv. — Gehen wir auf
dem Kreise z. B. in der Richtung weiter, in der \(U\) positiv wird, so wird in
dem nun folgenden Punkte, wo \(T=0\) ist, \(\mathrm{X}\) positiv imaginär, und wenn man
immer in demselben Sinne weiter geht, in dem nächsten Punkte, wo wieder
\(U=0\) ist, negativ reell, in dem nächsten Punkte, wo \(T=0\) ist, negativ
imaginär und in dem darauf folgenden Punkte, wo \(U=0\) ist, wieder reell
positiv; das Argument von \(\mathrm{X}\), das heisst also \(\operatorname{arctg}U/T\), hat sich demnach
um \(2\pi\) vermehrt. Wenn wir nun die ganze Kreisperipherie einmal durchlaufen,
so spielt sich das eben Geschilderte genau \(n\)-mal ab, so dass \(\arg\mathrm{X}\)\par
\clearpage\SourceLeaf{147}\ArithmeticHead{62}{L}{L. SCHLESINGER, ÜBER GAUSS’ FUNKTIONENTHEORETISCHE ARBEITEN.}
\noindent den Zuwachs \(2n\pi\) erfahren wird\footnote{Vergl. \textsc{A. Fraenkel}, Jahresbericht der D. M.-V. 31, 1922, S. 234 ff.}, denselben Zuwachs also, den auch das
Argument \(n\varphi\) von \(x^n\) erfährt, wie es ja sein muss, weil \(\arg(1+A/x+\ldots)\) doch
ungeändert bleibt. Bei einmaligem Durchlaufen des grossen Kreises erfährt
folglich \(\log\mathrm{X}\) die Änderung \(2n\pi i\), und nach den uns geläufigen funktionentheoretischen
Prinzipien folgt hieraus ohne weiteres, dass innerhalb des Kreises
genau \(n\) nach ihrer Vielfachheit gezählte Nullstellen von \(\mathrm{X}\) liegen müssen. —
Dass \textsc{Gauss} im Jahre 1799 nicht im Besitz dieser Prinzipien war, ist sicher,
wir werden im Laufe unserer Darstellung sehen, wie weit sich dieselben später
bei ihm entwickelt haben.
